% !TEX program = xelatex
\documentclass[12pt,a4paper]{article}
\usepackage[T1]{fontenc}
%\usepackage{tgpagella, eulervm}

% --- XeLatex/LuaLatex ---
\usepackage[no-math]{fontspec}
\IfFontExistsTF{Fira Sans}{
	\setmainfont{Fira Sans}
}

% % --- This is for pdflatex ---
%\usepackage[light, sfdefault]{FiraSans}

\usepackage[italic]{mathastext}
\usepackage[bahasa]{babel}
\usepackage{booktabs}
\usepackage{geometry}
\usepackage{tabularx}
\usepackage{graphicx}
\usepackage{float}
\usepackage{xcolor}
\usepackage{hyperref}
\hypersetup{
	colorlinks=true,  % Colours links instead of boxes
	urlcolor=blue,    % Colour for external hyperlinks
	linkcolor=blue,   % Colour of internal links (sections, etc.)
	citecolor=red,    % Colour of citations
	bookmarks=true,   % Adds PDF bookmarks
	hypertexnames=true
}

\geometry{left=2cm, right=2cm, top=2cm}
\title{\textbf{Panduan Pengamatan Okultasi HIP 35933 oleh Asteroid 1201 Strenua}}
\author{Agus T. P. Jatmiko}
\hyphenation{me-la-ku-kan peng-amat-an deng-an}
\begin{document}
\maketitle
\noindent\rule{\linewidth}{0.4pt}

\subsection*{Pengantar}
Dokumen ini berisi detail terkait pengambilan data okultasi, baik pada waktu latihan maupun waktu pengamatan sebenarnya.
\subsection*{Langkah 1: Persiapan awal}
\begin{enumerate}
	\item Lakukan persiapan setidaknya \textbf{3 jam} sebelum waktu \textit{event}.
	\item Jika menggunakan NTP timing, pastikan PC/laptop yang digunakan tersinkronisasi dengan server jam atom melalui jaringan internet setidaknya \textbf{2 jam} sebelum waktu pengamatan.
	
	Jika menggunakan GPS, pastikan GPS sudah dalam keadaan \textit{locked} sebelum melakukan pengambilan data.
	\item Pasang instrumen selayaknya pengamatan astronomi secara umum. Lakukan \textit{polar alignment} dan perbaiki \textit{pointing model} jika diperlukan.
	\item Ambil data kalibrasi FLAT dengan \textit{setting} instrumen yang sama (opsional, tapi direkomendasikan).  
\end{enumerate}
\subsection*{Langkah 2: \textit{Pointing} ke bintang target}
\begin{enumerate}
	\item Detail bintang target bisa dilihat pada tabel berikut.
	\begin{table}[H]
		\centering
		\large
		\begin{tabular}{ll}
			\toprule
			Target &: HIP 35933, HD 58050  \\
			RA (J2000) &: $7^\mathrm{h} 24^\mathrm{m} 27.648^\mathrm{s}$  \\
			DEC (J2000) &: $15^\circ$ 31' 1.89''  \\
			Waktu \textit{event} &: 12:41.8 UT (lokasi: Obs. Bosscha) \\
			\bottomrule
		\end{tabular}
	\end{table}
	\item Setelah teleskop mengarah ke target, gunakan \textit{finding chart} pada Gambar~\ref{fig:findingchart} untuk memastikan bahwa teleskop sudah benar mengarah ke bintang target. (\underline{catatan:} untuk teleskop yang sudah terkalibrasi dengan baik (\textit{polar aligned}) dan memiliki akurasi pointing tinggi bisa melewati tahap ini).
	\item Pastikan bintang dalam keadaan fokus saat diamati dengan teleskop. Lakukan koreksi jika belum.
\end{enumerate}

\subsection*{Langkah 3: Pengambilan Data}
\begin{enumerate}
	\item Pastikan bintang target berada di tengah \textit{frame} sesaat sebelum pengambilan data. \textcolor{red}{Jangan menggeser bintang target saat pengambilan data meskipun bintang bergeser. Jika \textit{tracking} teleskop kurang baik, pastikan peletakan target dilakukan dengan pertimbangan bahwa target akan tetap di dalam \textit{frame} selama pengambilan data.}
	\item Pastikan \textit{setting} pada SharpCap sudah sesuai dengan Panduan \textit{Setting} Kamera (dokumen terpisah).
	\item Pengambilan data dilakukan \textcolor{blue}{$\sim 5$ menit sebelum waktu okultasi dan dihentikan $\sim 5$ menit setelah waktu okultasi}, dengan detail sebagai berikut.
	\begin{enumerate}
		\item Untuk \textbf{malam latihan}, ambil total 10 menit data dari pukul \textbf{12:36.8 UT (19:36.8 WIB)} s.d. \textbf{12:46.8 UT (19:46.8 WIB)}.
		
		 Gunakan kesempatan ini untuk mencoba-coba waktu paparan (dimulai dari nilai rekomendasi) untuk menghasilkan kualitas data dengan SNR yang cukup tanpa mengorbankan terlalu banyak titik sampel.
		\item Untuk \textbf{malam gladi}, ambil total 10 menit data dari pukul \textbf{12:36.8 UT (19:36.8 WIB)} s.d. \textbf{12:46.8 UT (19:46.8 WIB)}.
		
		Kesempatan ini digunakan untuk finalisasi konfigurasi instrumen dan \textit{setting} kamera. Konfigurasi terbaik dari sesi ini akan dipertahankan untuk digunakan pada malam \textit{event}.
		\item Untuk \textbf{malam \textit{event}}, ambil total 10 menit data dari pukul \textbf{12:36.8 UT (19:36.8 WIB)} s.d. \textbf{12:46.8 UT (19:46.8 WIB)}. 
	\end{enumerate}
	\underline{Catatan umum}: 
	\begin{itemize}
		\item Sesuaikan waktu lokal/waktu zona pengamatan dengan lokasi \textit{chord} masing--masing.
		\item Pantau apakah ada \textit{drop frames}. Catat pada log--sheet jika ada (dan berapa banyak). Catat pula waktu awal dan akhir pengambilan data pada log--sheet.
		\item Gunakan log--sheet yang berbeda untuk setiap malam pengamatan. Beri judul yang sesuai (misal malam latihan, malam gladi, malam \textit{event})
	\end{itemize}
	\underline{Catatan log--sheet}:
	\begin{itemize}
		\item Gunakan/tulis waktu pengamatan dalam UT.
		\item Tulis lokasi pengamatan (bujur, lintang, ketinggian).
		\item Tulis konfigurasi sistem yang digunakan (teleskop, kamera, mode \textit{timing} (GPS/NTP), sistem operasi).
		\item Tulis nama lengkap semua orang yang terlibat dalam tim.
		\item Tulis inisial/singkatan/\textit{nick name} dari lokasi pengamatan dengan huruf kapital dan terdiri dari satu kata saja (contoh: Observatorium Bosscha: BOSSCHA).
	\end{itemize}
%	\pagebreak
	\textbf{DO \& DON'T}
	
	\underline{DO}:
	\begin{itemize}
		\item Kontrol pencahayaan. Lindungi teleskop dari paparan cahaya lingkungan (layar laptop, senter, dll)
		\item Jika memungkinkan, usahakan ada bintang lain di sekitar bintang target yang ikut terekam dalam satu \textit{frame}. Bintang ini akan berguna saat ekstraksi data. Bintang target tidak perlu terlalu sempurna berada di tengah \textit{frame}.
	\end{itemize}
	\underline{DON'T}
	\begin{itemize}
		\item Berjalan di sekitar teleskop saat pengambilan data.
		\item Menyalakan senter/lampu saat pengambilan data.
		\item Menggerakkan teleskop untuk \textit{guiding} bintang target saat pengambilan data. 		
	\end{itemize} 
\end{enumerate}
% TODO: \usepackage{graphicx} required
\begin{figure}
	\centering
	\includegraphics[width=0.65\linewidth]{finding_chart}
	\caption{\textit{Finding chart} untuk bintang target dengan FOV $40' \times 40'$.}
	\label{fig:findingchart}
\end{figure}

\subsection*{Langkah 4: Setelah pengamatan}
\begin{enumerate}
	\item Pastikan log--sheet sudah dilengkapi. Tulis semua informasi yang relevan terkait kondisi pengamatan.
	\item Unggah data pengamatan dan FLAT (jika ada) ke tautan Google Drive yang disediakan sesegera mungkin. Masukkan file--file tersebut ke dalam folder dengan format nama folder \textcolor{blue}{\textsc{\textbf{[nama chord]}}} (contoh: \textcolor{blue}{\textbf{BOSSCHA}}) dan ubah nama filenya dengan format: 
	
	\textcolor{blue}{
	\textsc{\textbf{[jenis pengamatan (latihan/gladi/event)]\_[nama chord]\_[tahun bulan tanggal]\_[tipe data (target/flat)]}}.adv
	}
	
	Contoh: \textcolor{blue}{\textbf{GLADI\_BOSSCHA\_20260401\_TARGET}.adv} 
%	
%	dan bagikan tautannya ke koordinator kegiatan.
	
	\underline{Catatan:}
	\begin{itemize}
%		\item Jika ada \textit{chord}/titik pengamatan yang keberatan membagikan data mentahnya secara langsung, PI \textit{chord} tersebut bertanggung jawab untuk mengolah data tersebut menjadi kurva cahaya dan menyimpannya ke dalam format .csv dengan penamaan file seperti tersebut di atas. Pastikan kalibrasi telah dilakukan (jika ada) sebelum mengekstrak kurva cahayanya. Unggah hasil ekstraksi tersebut ke tautan yang disediakan.
		\item Jika format file yang digunakan adalah .FITS, letakkan file--file ini ke dalam sub--folder di bawah folder utama setiap \textit{chord} dengan penamaan sesuai Langkah 4 Poin 2 di atas (contoh: \textcolor{blue}{\textbf{BOSSCHA/GLADI\_BOSSCHA\_20260401\_TARGET/}$\cdots$.fits})
	\end{itemize} 
\end{enumerate}
 
\end{document}